\section{Linear algebra and numerics}

\entry{Gauss mod p}{$O(n^2 m)$ \quad mem $nm$}{%
Needs \code{power}. Solves $Ax \equiv b \pmod p$, $p$ prime, $A$ is $n \times m$ ($n$
equations, $m$ unknowns). Returns \code{-1} if there is no solution, else the rank, and
\code{x} is one solution with free variables 0. Rank $< m$ means $p^{m - \text{rank}}$
solutions in total --- that is how ``count the solutions'' questions are answered.
Exact, so no \code{eps}. Ruiming's \code{mod\_elimination}, rewritten 0-indexed with
the inverse taken by \code{power}.}

%DEP modbasics
%LABEL gaussMod
\begin{lstlisting}
int gaussMod(vector<vi> a, vi b, vi& x, ll p) {
    int n = (int)a.size(), m = (int)a[0].size(), rank = 0;
    for (auto& r : a) for (ll& v : r) v = (v % p + p) % p;
    for (ll& v : b) v = (v % p + p) % p;
    vector<int> where(m, -1);
    for (int col = 0; col < m && rank < n; col++) {
        int sel = -1;
        for (int i = rank; i < n; i++)
            if (a[i][col]) { sel = i; break; }
        if (sel < 0) continue;
        swap(a[sel], a[rank]); swap(b[sel], b[rank]);
        where[col] = rank;
        ll iv = power(a[rank][col], p - 2, p);
        for (int j = col; j < m; j++)
            a[rank][j] = a[rank][j] * iv % p;
        b[rank] = b[rank] * iv % p;
        for (int i = 0; i < n; i++)
            if (i != rank && a[i][col]) {
                ll c = a[i][col];
                for (int j = col; j < m; j++)
                    a[i][j] = ((a[i][j] - c * a[rank][j]) % p
                               + p) % p;
                b[i] = ((b[i] - c * b[rank]) % p + p) % p;
            }
        rank++;
    }
    for (int i = rank; i < n; i++) if (b[i]) return -1;
    x.assign(m, 0);
    for (int j = 0; j < m; j++)
        if (where[j] >= 0) x[j] = b[where[j]];
    return rank;
}
\end{lstlisting}

\entry{Gauss mod 2 (XOR equations)}{$O(n^2 m / 64)$}{%
The same thing over bits: each equation is ``the XOR of some variables equals 0 or 1''.
Rows are \code{bitset}s, so $n = m = 2000$ is instant. Lights-out puzzles, ``flip
switches so every lamp is off'', parity constraints. Fill \code{A[i][j] = 1} when
variable \code{j} appears in equation \code{i}, and \code{b[i]} is the right side.
Returns \code{-1} if impossible, else the rank; $2^{m - \text{rank}}$ solutions exist.
Destroys \code{A} and \code{b}. Bits at columns $\ge m$ must be 0. KACTL's, replacing
Ruiming's \code{xor\_gauss} (same algorithm, general interface).}

%LABEL gaussXor
\begin{lstlisting}
typedef bitset<2048> bs;
int gaussXor(vector<bs>& A, vector<int>& b, bs& x, int m) {
    int n = (int)A.size(), rank = 0, br;
    vector<int> col(m); iota(all(col), 0);
    rep(i, 0, n) {
        for (br = i; br < n; ++br) if (A[br].any()) break;
        if (br == n) {
            rep(j, i, n) if (b[j]) return -1;
            break;
        }
        int bc = (int)A[br]._Find_next(i - 1);
        swap(A[i], A[br]); swap(b[i], b[br]);
        swap(col[i], col[bc]);
        rep(j, 0, n) if (A[j][i] != A[j][bc]) {
            A[j].flip(i); A[j].flip(bc);
        }
        rep(j, i + 1, n) if (A[j][i]) {
            b[j] ^= b[i]; A[j] ^= A[i];
        }
        rank++;
    }
    x = bs();
    for (int i = rank; i--;) {
        if (!b[i]) continue;
        x[col[i]] = 1;
        rep(j, 0, i) b[j] ^= A[j][i];
    }
    return rank;
}
\end{lstlisting}

\entry{XOR basis}{add $O(\log C)$ \quad mem 63 words}{%
Keeps a small set of numbers (at most 63) whose XOR combinations are exactly the XORs
of all subsets of everything added so far. Answers: ``can $x$ be made as a XOR of some
of the numbers?'' (\code{can}), ``largest XOR of a subset'' (\code{maxXor()}), ``largest
\code{x \^{} subset}'' (\code{maxXor(x)}), ``$k$-th smallest reachable value''
(\code{kth}, counting 0 as the 0-th).\\[0.3ex]
\code{add} returns \code{false} when the number was already reachable: then the empty
subset is not the only way to get 0, which matters if the problem forbids empty subsets
(if \emph{every} add returned \code{true}, 0 is not reachable by a non-empty subset, so
shift \code{k} by one). Count of distinct reachable values: $2^{\text{rank}}$. Works
for numbers $< 2^{63}$; combine with a DFS for ``max XOR path in a graph'' (add every
cycle's XOR). Ruiming's \code{xor\_basis}, as a struct.}

%LABEL xorbasis
\begin{lstlisting}
struct XorBasis {
    static const int B = 63;
    ll b[B] = {}; int rank = 0;
    bool add(ll x) {
        for (int i = B - 1; i >= 0; i--) if (x >> i & 1) {
            if (!b[i]) { b[i] = x; rank++; return true; }
            x ^= b[i];
        }
        return false;                    // already reachable
    }
    bool can(ll x) {
        for (int i = B - 1; i >= 0; i--) if (x >> i & 1) {
            if (!b[i]) return false;
            x ^= b[i];
        }
        return true;
    }
    ll maxXor(ll x = 0) {
        for (int i = B - 1; i >= 0; i--) x = max(x, x ^ b[i]);
        return x;
    }
    ll kth(ll k) {                // 0-th = 0; -1 if too big
        vector<ll> v;
        for (int i = 0; i < B; i++) if (b[i]) {
            for (int j = i - 1; j >= 0; j--)  // reduce
                if (b[i] >> j & 1) b[i] ^= b[j];
            v.push_back(b[i]);
        }
        if (v.size() < 63 && (k >> v.size())) return -1;
        ll r = 0;
        rep(i, 0, (ll)v.size()) if (k >> i & 1) r ^= v[i];
        return r;
    }
};
\end{lstlisting}

\entry{Determinant}{$O(n^3)$}{%
\code{det} for real matrices (partial pivoting); \code{detMod} is exact modulo
\emph{any} \code{md}, prime or not, because it only uses Euclid-style row subtraction
(no division). Reduce the input mod \code{md} first.\\[0.3ex]
\textbf{Count spanning trees} (Kirchhoff): build the Laplacian $L$ ($L_{ii}$ = degree,
$L_{ij} = -$(number of edges $i$--$j$)), delete row 0 and column 0, take the
determinant. Directed, trees pointing into root $r$: $L_{ii}$ = out-degree, delete row
and column $r$. Also: $\det = 0$ iff the rows are dependent. From KACTL.}

%LABEL det
\begin{lstlisting}
ld det(vector<vector<ld>> a) {
    int n = (int)a.size(); ld res = 1;
    rep(i, 0, n) {
        int b = i;
        rep(j, i + 1, n)
            if (fabsl(a[j][i]) > fabsl(a[b][i])) b = j;
        if (i != b) swap(a[i], a[b]), res *= -1;
        res *= a[i][i];
        if (res == 0) return 0;
        rep(j, i + 1, n) {
            ld v = a[j][i] / a[i][i];
            if (v != 0) rep(k, i + 1, n) a[j][k] -= v*a[i][k];
        }
    }
    return res;
}
ll detMod(vector<vi> a, ll md) {
    int n = (int)a.size(); ll ans = 1;
    rep(i, 0, n) {
        rep(j, i + 1, n) while (a[j][i] != 0) {   // gcd step
            ll t = a[i][i] / a[j][i];
            if (t) rep(k, i, n)
                a[i][k] = (a[i][k] - a[j][k] * t) % md;
            swap(a[i], a[j]);
            ans *= -1;
        }
        ans = ans * a[i][i] % md;
        if (!ans) return 0;
    }
    return (ans + md) % md;
}
\end{lstlisting}

\entry{Berlekamp-Massey + linear recurrence}{$O(n^2)$, then $O(n^2 \log k)$}{%
Needs \code{power}. \textbf{The ``guess the formula'' tool.} If the answer for $n$
follows \emph{some} linear recurrence $a_i = c_1 a_{i-1} + \dots + c_d a_{i-d}$ (mod a
prime) --- which is true for every ``count paths / tilings / strings with a fixed
small automaton'' problem and anything a matrix power could solve --- then:\\[0.3ex]
1. Brute force the first $2d + 10$ terms (more is safer), mod the prime.\\
2. \code{tr = berlekampMassey(terms, md)} returns $c_1, \dots, c_d$.\\
3. \code{linearRec(terms, tr, k, md)} gives term $k$ (0-indexed), $k$ up to $10^{18}$.\\[0.3ex]
Much faster and simpler than building the matrix by hand: you never have to find the
recurrence yourself. If \code{tr} comes out as long as half your input, you did not
give enough terms or there is no linear recurrence. Example:
\code{berlekampMassey(\{0,1,1,3,5,11\}, md)} is \code{\{1, 2\}}. From KACTL.}

%DEP modbasics
%LABEL bm
\begin{lstlisting}
vi berlekampMassey(vi s, ll md) {
    int n = (int)s.size(), L = 0, m = 0;
    vi C(n), B(n), T;
    C[0] = B[0] = 1;
    ll b = 1;
    rep(i, 0, n) { ++m;
        ll d = s[i] % md;
        rep(j, 1, L + 1) d = (d + C[j] * s[i - j]) % md;
        if (!d) continue;
        T = C; ll coef = d * power(b, md - 2, md) % md;
        rep(j, m, n) C[j] = (C[j] - coef * B[j - m]) % md;
        if (2 * L > i) continue;
        L = i + 1 - L; B = T; b = d; m = 0;
    }
    C.resize(L + 1); C.erase(C.begin());
    for (ll& x : C) x = (md - x) % md;
    return C;
}
ll linearRec(vi S, vi tr, ll k, ll md) {   // S[k]
    int n = (int)tr.size();
    auto combine = [&](vi a, vi b) {
        vi res(n * 2 + 1);
        rep(i, 0, n + 1) rep(j, 0, n + 1)
            res[i + j] = (res[i + j] + a[i] * b[j]) % md;
        for (int i = 2 * n; i > n; --i) rep(j, 0, n)
            res[i-1-j] = (res[i-1-j] + res[i] * tr[j]) % md;
        res.resize(n + 1);
        return res;
    };
    vi pol(n + 1), e(pol);
    pol[0] = e[1] = 1;
    for (++k; k; k /= 2) {
        if (k % 2) pol = combine(pol, e);
        e = combine(e, e);
    }
    ll res = 0;
    rep(i, 0, n) res = (res + pol[i + 1] * S[i]) % md;
    return res;
}
\end{lstlisting}

\entry{Linear programming (simplex)}{fast in practice}{%
Maximise $c^T x$ subject to $Ax \le b$ and $x \ge 0$. Returns the best value, or
\code{-inf} if the rules contradict each other, or \code{inf} if it can grow forever;
\code{x} gets an optimal point.\\[0.3ex]
Turning a problem into this form: $\ge$ constraints are $\le$ with both sides negated;
an equality is one $\le$ and one $\ge$; minimise $f$ = maximise $-f$; a variable that
may be negative is $x = x^+ - x^-$ with two non-negative variables. Use it when the
problem is ``choose real amounts, every rule is a linear inequality, optimise a linear
total'' --- e.g. diet / mixing / scheduling with fractional amounts. Answers are
\code{double}: print with enough digits. Integer answers are a different (hard)
problem. From KACTL.\\[0.3ex]
\code{LPSolver::vd b = \{1, 1, -4\}, c = \{-1, -1\}, x;}\\
\code{LPSolver(A, b, c).solve(x);} with \code{A = \{\{1,-1\},\{-1,1\},\{-1,-2\}\}}.}

%LABEL simplex
\begin{lstlisting}
#define MP make_pair
#define ltj(X) if (s==-1 || MP(X[j],N[j]) < MP(X[s],N[s])) s=j
struct LPSolver {
    using T = double; using vd = vector<T>;
    static constexpr T eps = 1e-8;
    static constexpr T inf = numeric_limits<T>::infinity();
    int m, n; vector<int> N, B; vector<vd> D;
    LPSolver(const vector<vd>& A, const vd& b, const vd& c)
        : m((int)b.size()), n((int)c.size()), N(n + 1),
          B(m), D(m + 2, vd(n + 2)) {
        rep(i, 0, m) rep(j, 0, n) D[i][j] = A[i][j];
        rep(i, 0, m) {
            B[i] = n + i; D[i][n] = -1; D[i][n+1] = b[i];
        }
        rep(j, 0, n) { N[j] = j; D[m][j] = -c[j]; }
        N[n] = -1; D[m+1][n] = 1;
    }
    void pivot(int r, int s) {
        T *a = D[r].data(), inv = 1 / a[s];
        rep(i, 0, m + 2) if (i != r && abs(D[i][s]) > eps) {
            T *b = D[i].data(), inv2 = b[s] * inv;
            rep(j, 0, n + 2) b[j] -= a[j] * inv2;
            b[s] = a[s] * inv2;
        }
        rep(j, 0, n + 2) if (j != s) D[r][j] *= inv;
        rep(i, 0, m + 2) if (i != r) D[i][s] *= -inv;
        D[r][s] = inv;
        swap(B[r], N[s]);
    }
    bool simplex(int phase) {
        int x = m + phase - 1;
        for (;;) {
            int s = -1;
            rep(j, 0, n + 1) if (N[j] != -phase) ltj(D[x]);
            if (D[x][s] >= -eps) return true;
            int r = -1;
            rep(i, 0, m) {
                if (D[i][s] <= eps) continue;
                if (r == -1 || MP(D[i][n+1] / D[i][s], B[i])
                             < MP(D[r][n+1] / D[r][s], B[r]))
                    r = i;
            }
            if (r == -1) return false;
            pivot(r, s);
        }
    }
    T solve(vd& x) {
        int r = 0;
        rep(i, 1, m) if (D[i][n+1] < D[r][n+1]) r = i;
        if (D[r][n+1] < -eps) {
            pivot(r, n);
            if (!simplex(2) || D[m+1][n+1] < -eps)
                return -inf;
            rep(i, 0, m) if (B[i] == -1) {
                int s = 0;
                rep(j, 1, n + 1) ltj(D[i]);
                pivot(i, s);
            }
        }
        bool ok = simplex(1); x = vd(n);
        rep(i, 0, m) if (B[i] < n) x[B[i]] = D[i][n+1];
        return ok ? D[m][n+1] : inf;
    }
};
\end{lstlisting}

\entry{AND / OR / XOR convolution}{$O(n \log n)$, $n = 2^k$}{%
\code{c[z]} = sum of \code{a[x]*b[y]} over all pairs with \code{x OP y == z}, where OP
is AND (\code{mode 0}), OR (\code{1}) or XOR (\code{2}). The same job as \code{conv} for
$+$, but on bit masks. Pad both arrays with zeros to the same power of two.\\[0.3ex]
Uses: ``number of pairs whose XOR is $z$'' for all $z$ at once; combining two
distributions over subsets; Nim-like games (XOR of piles) summed over choices.
\code{FST} on its own: OR forward = sum over subsets (same as \code{sos}), AND forward
= sum over supersets. Values grow like $n \cdot \max^2$; to work mod a prime, add
\code{\% md} to every line and replace the final \code{/ n} by multiplying with
\code{power(n, md-2, md)}. b5xu's \code{bitwise-xor-convolution}, fixed (the original
does not compile) and extended with KACTL's AND/OR.}

%LABEL fst
\begin{lstlisting}
void FST(vi& a, bool inv, int mode) {    // 0 AND 1 OR 2 XOR
    int n = (int)a.size();
    for (int s = 1; s < n; s *= 2)
        for (int i = 0; i < n; i += 2 * s) rep(j, i, i + s) {
            ll &u = a[j], &v = a[j + s];
            if (mode == 0) u += inv ? -v : v;
            else if (mode == 1) v += inv ? -u : u;
            else tie(u, v) = pii(u + v, u - v);
        }
    if (inv && mode == 2) for (ll& x : a) x /= n;
}
vi bitConv(vi a, vi b, int mode) {     // same size, 2^k
    FST(a, 0, mode); FST(b, 0, mode);
    rep(i, 0, (ll)a.size()) a[i] *= b[i];
    FST(a, 1, mode);
    return a;
}
\end{lstlisting}

\entry{Polynomial interpolation}{$O(n^2)$ / $O(n \log p)$}{%
Any $n$ points with different $x$ lie on exactly one polynomial of degree $< n$.
\code{interpolate} returns its coefficients (\code{double}; inaccurate for large $n$).
\code{lagrange} is the exact, more useful one: given $y_i = p(i)$ for
$i = 0..n-1$ mod a prime $\ge n$, it returns $p(X)$ for a huge $X$.\\[0.3ex]
Typical use: ``the answer is a polynomial in $n$ of degree $d$'' --- e.g.
$\sum_{i \le n} i^k$ has degree $k+1$, and many counting answers over a growing box are
polynomials. Brute force $d+1$ (or a few more) small values, then \code{lagrange} at
the real $n$. Linear recurrence instead of polynomial? use Berlekamp-Massey.}

%DEP modbasics
%LABEL interp
\begin{lstlisting}
vector<ld> interpolate(vector<ld> x, vector<ld> y) {
    int n = (int)x.size();
    vector<ld> res(n), temp(n);
    rep(k, 0, n - 1) rep(i, k + 1, n)
        y[i] = (y[i] - y[k]) / (x[i] - x[k]);
    ld last = 0; temp[0] = 1;
    rep(k, 0, n) rep(i, 0, n) {
        res[i] += y[k] * temp[i];
        swap(last, temp[i]);
        temp[i] -= last * x[k];
    }
    return res;                      // res[i] * x^i
}
ll lagrange(const vi& y, ll X, ll md) {  // y[i] = p(i)
    int n = (int)y.size(); X %= md;
    vi pre(n + 1, 1), suf(n + 1, 1), f(n, 1);
    rep(i, 0, n) pre[i+1] = pre[i] * ((X - i + md) % md) % md;
    down(i, n, 0) suf[i] = suf[i+1] * ((X - i + md) % md)%md;
    rep(i, 1, n) f[i] = f[i-1] * i % md;
    ll res = 0;
    rep(i, 0, n) {
        ll t = y[i] % md * (pre[i] * suf[i+1] % md) % md
             * power(f[i] * f[n-1-i] % md, md - 2, md) % md;
        res = ((n - 1 - i) & 1) ? (res - t + md) % md
                                : (res + t) % md;
    }
    return res;
}
\end{lstlisting}

\entry{Numerical integration}{adaptive Simpson}{%
\code{quad(a, b, f)} integrates \code{f} from \code{a} to \code{b}. It splits the
interval in halves only where the function is hard, so it is fast for smooth functions
and still fine for ones with a kink. Area of a weird shape: integrate its height
(e.g. the length of a vertical cut through a union of circles) over $x$. Volume:
nest it, \code{quad(.., [\&](ld x)\{ return quad(.., [\&](ld y)\{ ... \}); \})}.\\[0.3ex]
Trouble with jumps (a function that is 0 almost everywhere): split at the known jump
points yourself and call it per piece. From KACTL.}

%LABEL quad
\begin{lstlisting}
template<class F> ld simp(F& f, ld a, ld b) {
    return (f(a) + 4 * f((a + b) / 2) + f(b)) * (b - a) / 6;
}
template<class F>
ld simpRec(F& f, ld a, ld b, ld eps, ld S) {
    ld c = (a + b) / 2;
    ld S1 = simp(f, a, c), S2 = simp(f, c, b), T = S1 + S2;
    if (fabsl(T - S) <= 15 * eps || b - a < 1e-10)
        return T + (T - S) / 15;
    return simpRec(f, a, c, eps / 2, S1)
         + simpRec(f, c, b, eps / 2, S2);
}
template<class F> ld quad(ld a, ld b, F f, ld eps = 1e-8) {
    return simpRec(f, a, b, eps, simp(f, a, b));
}
\end{lstlisting}
